% GENERATED FILE. Edit canonical lessons, then run npm run generate:books.
% canonical-source-hash: fnv1a64:de23b61826ba53ed
% canonical-lessons: RU-C383-postuplenie, RU-C383-abiturient, RU-C383-vstupitelnyy, RU-C383-fakultet, RU-C383-stipendiya

\chapter{Getting into university}
\label{ch:ru-getting-into-university}
\hlchaptermodality{383}

\begin{chapteropening}
\textbf{By the end of this chapter:} \emph{I can talk about getting into university: applicants, entrance exams, the faculty and a grant.}

Complete the last lesson of chapter 383: I can talk about getting into university: applicants, entrance exams, the faculty and a grant.
\end{chapteropening}

\section[\texorpdfstring{\ru{поступление} (postupléniye)}{поступление (postupléniye)}]{\ru{поступление} (postupléniye) — admission (to a university); entry}

\label{lesson:RU-C383-postuplenie}



\begin{quote}
Before the new one: say the Russian for preparation, then the Russian for to prepare oneself.
\end{quote}

\subsection*{You'll want to know: \ru{поступление}}
\textbf{\ru{поступление}} (\emph{postupléniye}) — \textquotedblleft{}admission (to a university); entry\textquotedblright{}.

\begin{quote}
\ru{Поступление} \ru{в} \ru{университет} — \ru{важный} \ru{шаг}.
\end{quote}

\begin{quote}
\emph{Postupléniye v universitét — vázhnyy shag.}
\end{quote}

\begin{quote}
Getting into university is an important step.
\end{quote}

Neuter. From \emph{\ru{поступать}}, to enter a university. It also means money coming in: \emph{\ru{поступление} \ru{денег}}.

\begin{cousinweb}
\textbf{\ru{поступление}} is built on \emph{\ru{поступить}} \textquotedblleft{}to enter\textquotedblright{}, from \emph{\ru{по}-} and \emph{\ru{ступить}} \textquotedblleft{}to step\textquotedblright{}, an inherited Slavic verb: the step into a university.
\end{cousinweb}

\begin{grammarlens}[title={What you've built}]
One more word for telling, explaining and describing at B1.
\end{grammarlens}

\subsection*{Your turn}
\begin{itemize}
\raggedright
  \item \emph{Say it:} \emph{postupléniye}
  \item \emph{Say it:} \emph{postupléniye}, once more
  \item \emph{Say it:} a sentence of your own with \emph{postupléniye}
  \item \emph{Recall:} say the Russian for preparation, then the Russian for to prepare oneself, then say \emph{postupléniye} again
\end{itemize}

\begin{tcolorbox}[breakable,colback=teal!4,colframe=teal!35!black,title={Before you move on}]
What does \ru{поступление} mean? (\textquotedblleft{}Admission (to a university); entry\textquotedblright{}.) Say it once more.
\end{tcolorbox}
\section[\texorpdfstring{\ru{абитуриент} (abituriyént)}{абитуриент (abituriyént)}]{\ru{абитуриент} (abituriyént) — a university applicant}

\label{lesson:RU-C383-abiturient}



\begin{quote}
Before the new one: say the Russian for to prepare oneself, then the Russian for admission.
\end{quote}

\subsection*{You'll want to know: \ru{абитуриент}}
\textbf{\ru{абитуриент}} (\emph{abituriyént}) — \textquotedblleft{}a university applicant\textquotedblright{}.

\begin{quote}
\ru{Абитуриенты} \ru{сдают} \ru{экзамены} \ru{летом}.
\end{quote}

\begin{quote}
\emph{Abituriyénty sdayút ekzámeny létom.}
\end{quote}

\begin{quote}
University applicants sit their exams in the summer.
\end{quote}

Masculine. Someone applying to a university; in German the same word means a school leaver.

\begin{cousinweb}
\textbf{\ru{абитуриент}} was borrowed from German \emph{Abiturient} (EtymDB), from Latin \emph{abiturīre} \textquotedblleft{}to wish to leave\textquotedblright{}, built on \emph{abīre} \textquotedblleft{}to go away\textquotedblright{}, from \emph{īre} \textquotedblleft{}to go\textquotedblright{}: a pupil about to leave school.
\end{cousinweb}

\begin{grammarlens}[title={What you've built}]
One more word for telling, explaining and describing at B1.
\end{grammarlens}

\subsection*{Your turn}
\begin{itemize}
\raggedright
  \item \emph{Say it:} \emph{abituriyént}
  \item \emph{Say it:} \emph{abituriyént}, once more
  \item \emph{Say it:} a sentence of your own with \emph{abituriyént}
  \item \emph{Recall:} say the Russian for to prepare oneself, then the Russian for admission, then say \emph{abituriyént} again
\end{itemize}

\begin{tcolorbox}[breakable,colback=teal!4,colframe=teal!35!black,title={Before you move on}]
What does \ru{абитуриент} mean? (\textquotedblleft{}A university applicant\textquotedblright{}.) Say it once more.
\end{tcolorbox}
\section[\texorpdfstring{\ru{вступительный} (vstupítelnyy)}{вступительный (vstupítelnyy)}]{\ru{вступительный} (vstupítelnyy) — entrance (exam), introductory}

\label{lesson:RU-C383-vstupitelnyy}



\begin{quote}
Before the new one: say the Russian for admission, then the Russian for a university applicant.
\end{quote}

\subsection*{You'll want to know: \ru{вступительный}}
\textbf{\ru{вступительный}} (\emph{vstupítelnyy}) — \textquotedblleft{}entrance (exam), introductory\textquotedblright{}.

\begin{quote}
\ru{Вступительный} \ru{экзамен} \ru{был} \ru{трудным}.
\end{quote}

\begin{quote}
\emph{Vstupítelnyy ekzámen byl trúdnym.}
\end{quote}

\begin{quote}
The entrance exam was hard.
\end{quote}

\emph{\ru{Вступительные} \ru{экзамены}} — entrance exams. \emph{\ru{Вступительное} \ru{слово}} — an opening speech.

\begin{cousinweb}
\textbf{\ru{вступительный}} is built on \emph{\ru{вступить}} \textquotedblleft{}to enter\textquotedblright{}, from \emph{\ru{в}-} \textquotedblleft{}in\textquotedblright{} and \emph{\ru{ступить}} \textquotedblleft{}to step\textquotedblright{}, an inherited Slavic verb.
\end{cousinweb}

\begin{grammarlens}[title={What you've built}]
One more word for telling, explaining and describing at B1.
\end{grammarlens}

\subsection*{Your turn}
\begin{itemize}
\raggedright
  \item \emph{Say it:} \emph{vstupítelnyy}
  \item \emph{Say it:} \emph{vstupítelnyy}, once more
  \item \emph{Say it:} a sentence of your own with \emph{vstupítelnyy}
  \item \emph{Recall:} say the Russian for admission, then the Russian for a university applicant, then say \emph{vstupítelnyy} again
\end{itemize}

\begin{tcolorbox}[breakable,colback=teal!4,colframe=teal!35!black,title={Before you move on}]
What does \ru{вступительный} mean? (\textquotedblleft{}Entrance (exam), introductory\textquotedblright{}.) Say it once more.
\end{tcolorbox}
\section[\texorpdfstring{\ru{факультет} (fakultét)}{факультет (fakultét)}]{\ru{факультет} (fakultét) — a faculty (of a university)}

\label{lesson:RU-C383-fakultet}



\begin{quote}
Before the new one: say the Russian for a university applicant, then the Russian for entrance.
\end{quote}

\subsection*{You'll want to know: \ru{факультет}}
\textbf{\ru{факультет}} (\emph{fakultét}) — \textquotedblleft{}a faculty (of a university)\textquotedblright{}.

\begin{quote}
\ru{Сестра} \ru{учится} \ru{на} \ru{историческом} \ru{факультете}.
\end{quote}

\begin{quote}
\emph{Sestrá úchitsya na istorícheskom fakultéte.}
\end{quote}

\begin{quote}
My sister studies at the history faculty.
\end{quote}

Masculine. Students study \emph{\ru{на} \ru{факультете}}, with \emph{\ru{на}}.

\begin{cousinweb}
\textbf{\ru{факультет}} was borrowed from German \emph{Fakultät}, from Latin \emph{facultās} \textquotedblleft{}capability, ability, skill\textquotedblright{} (EtymDB): a branch of learning as a power of the mind.
\end{cousinweb}

\begin{grammarlens}[title={What you've built}]
One more word for telling, explaining and describing at B1.
\end{grammarlens}

\subsection*{Your turn}
\begin{itemize}
\raggedright
  \item \emph{Say it:} \emph{fakultét}
  \item \emph{Say it:} \emph{fakultét}, once more
  \item \emph{Say it:} a sentence of your own with \emph{fakultét}
  \item \emph{Recall:} say the Russian for a university applicant, then the Russian for entrance, then say \emph{fakultét} again
\end{itemize}

\begin{tcolorbox}[breakable,colback=teal!4,colframe=teal!35!black,title={Before you move on}]
What does \ru{факультет} mean? (\textquotedblleft{}A faculty (of a university)\textquotedblright{}.) Say it once more.
\end{tcolorbox}
\section[\texorpdfstring{\ru{стипендия} (stipéndiya)}{стипендия (stipéndiya)}]{\ru{стипендия} (stipéndiya) — a student grant, a scholarship}

\label{lesson:RU-C383-stipendiya}



\begin{quote}
Before the new one: say the Russian for entrance, then the Russian for a faculty.
\end{quote}

\subsection*{You'll want to know: \ru{стипендия}}
\textbf{\ru{стипендия}} (\emph{stipéndiya}) — \textquotedblleft{}a student grant, a scholarship\textquotedblright{}.

\begin{quote}
\ru{Он} \ru{получает} \ru{стипендию}, \ru{потому} \ru{что} \ru{хорошо} \ru{учится}.
\end{quote}

\begin{quote}
\emph{On polucháyet stipéndiyu, potomú chto khoroshó úchitsya.}
\end{quote}

\begin{quote}
He gets a grant because he does well at his studies.
\end{quote}

Feminine. A monthly payment to students.

\begin{cousinweb}
\textbf{\ru{стипендия}} is usually traced to Latin \emph{stīpendium} \textquotedblleft{}a tax, a soldier’s pay\textquotedblright{} (Lewis \& Short).
\end{cousinweb}

\begin{grammarlens}[title={What you've built}]
One more word for telling, explaining and describing at B1.
\end{grammarlens}

\subsection*{Your turn}
\begin{itemize}
\raggedright
  \item \emph{Say it:} \emph{stipéndiya}
  \item \emph{Say it:} \emph{stipéndiya}, once more
  \item \emph{Say it:} a sentence of your own with \emph{stipéndiya}
  \item \emph{Recall:} say the Russian for entrance, then the Russian for a faculty, then say \emph{stipéndiya} again
\end{itemize}

\begin{tcolorbox}[breakable,colback=teal!4,colframe=teal!35!black,title={Before you move on}]
What does \ru{стипендия} mean? (\textquotedblleft{}A student grant, a scholarship\textquotedblright{}.) Say it once more.
\end{tcolorbox}
